\section{Функция потерь и обучение}

\begin{equation}
\mathcal{L} = \sum_{m,e}
\Bigl[\lambda_{\text{крив}}\ell_{\text{крив}} + \lambda_{\text{tdi}}\ell_{\text{tdi}}
+ \lambda_{\max}\ell_{\max} + \lambda_{\text{скр}}\ell_{\text{скр}}\Bigr]
+ \lambda_{\text{полоса}}\ell_{\text{полоса}} + \lambda_{\text{st}}\ell_{\text{st}} + \Omega(W) .
\label{eq:loss}
\end{equation}

Первые четыре слагаемых --- правдоподобия четырёх источников из раздела~\ref{sec:gen}. Пятое обучает головы полос
$(\sigma^-,\sigma^+)$. Шестое --- клипованная L1 из \eqref{eq:strae}, посчитанная на обучающих полосах;
она выравнивает функцию потерь с метрикой соревнования, потому что правдоподобие оптимизирует не совсем то,
по чему считается лидерборд. Последнее --- регуляризатор \eqref{eq:omega}.

\subsection{Батчинг по молекулам, а не по строкам}

Единица батча --- молекула со всеми её наблюдениями по всем источникам и ферментам. Иначе теряется смысл конструкции: латентная $\theta_{m,e}$ общая для кривой, Emax и скрининга,
и градиенты от них должны сходиться на одном узле в пределах одного шага. Батч по строкам превращает конструкцию обратно
в обычную многозадачную модель с общим стволом.

\subsection{Двухэтапная калибровка}

Если учить $\psi_e = (s_e, b_e, \sigma_e)$ совместно с сетью с самого начала, есть вырожденное решение:
сеть подстраивает $\pi$ под удобную калибровку вместо обратного. Порядок такой.

\begin{enumerate}
\item На подмножестве, где есть и кривая, и скрининг, фиксируем $\pi$ равным опубликованному и подгоняем
12 скаляров $\psi$ напрямую. Задача низкоразмерная, решается за секунды; результат --- в
разделе~\ref{sec:calib}.
\item Обучаем сеть с замороженной $\psi$.
\item Размораживаем $\psi$ с маленьким шагом обучения на последних эпохах.
\end{enumerate}

\subsection{Веса источников}

Стартуем с $\lambda \propto 1/n_{\text{источника}}$, чтобы 17.5 тысяч точек скрининга не задавили
6.5 тысяч меток кривых. Потом пробуем гомоскедастичное взвешивание по неопределённости: каждому источнику
свой обучаемый $\log \sigma^2$, веса выводятся из него, и модель сама решает, какому источнику доверять.
Проверяем это абляцией, а не берём на веру: приём известен, но на маленьких выборках его собственный шум
иногда съедает пользу.
